% !TeX program = pdflatex
\documentclass[11pt,a4paper]{amsart}

\newcommand{\notelanguage}{finnish}
% Shared preamble for course notes and exercise sets.

\usepackage[T1]{fontenc}
\usepackage{lmodern}
\providecommand{\notelanguage}{english}
\usepackage[finnish,english]{babel}
\AtBeginDocument{\expandafter\selectlanguage\expandafter{\notelanguage}}

\newcommand{\notesfinnishlanguage}{finnish}
\ifx\notelanguage\notesfinnishlanguage
\newcommand{\notestheoremname}{Lause}
\newcommand{\noteslemmaname}{Lemma}
\newcommand{\notespropositionname}{Propositio}
\newcommand{\notescorollaryname}{Korollaari}
\newcommand{\notesdefinitionname}{Määritelmä}
\newcommand{\notesexamplename}{Esimerkki}
\newcommand{\notesproblemname}{Tehtävä}
\newcommand{\notesexercisename}{Tehtävä}
\newcommand{\notesremarkname}{Huomautus}
\newcommand{\notessolutionname}{Ratkaisu}
\else
\newcommand{\notestheoremname}{Theorem}
\newcommand{\noteslemmaname}{Lemma}
\newcommand{\notespropositionname}{Proposition}
\newcommand{\notescorollaryname}{Corollary}
\newcommand{\notesdefinitionname}{Definition}
\newcommand{\notesexamplename}{Example}
\newcommand{\notesproblemname}{Problem}
\newcommand{\notesexercisename}{Exercise}
\newcommand{\notesremarkname}{Remark}
\newcommand{\notessolutionname}{Solution}
\fi

\usepackage[a4paper,margin=30mm]{geometry}
\usepackage{microtype}
\usepackage{graphicx}
\usepackage{enumitem}

\newcommand{\exerciseheading}[2]{%
  \par\noindent
  {\large\bfseries #1\par}
  \vspace{0.1em}
  \noindent{\large #2\par}
  \vspace{1.5em}
}

\usepackage{amsmath}
\usepackage{amssymb}
\usepackage{amsthm}
\usepackage{mathtools}
\usepackage{aliascnt}

% Row operations beside the affected row, e.g. \rowop{+2R_1}.
\newcommand{\rowop}[1]{\mathrel{\big|}\,#1}

% Linear systems: the argument gives the number of alignment pairs.
\newenvironment{linearsystem}[1]
{\left\{
  \begin{alignedat}{#1}}
    {
  \end{alignedat}\right.}

  \usepackage[hidelinks,pdfusetitle]{hyperref}

  \numberwithin{equation}{section}
  \setlist{itemsep=0.25em,topsep=0.5em}

  \theoremstyle{plain}
  \newtheorem{theorem}{\notestheoremname}[section]
  \newaliascnt{lemma}{theorem}
  \newtheorem{lemma}[lemma]{\noteslemmaname}
  \aliascntresetthe{lemma}
  \newaliascnt{proposition}{theorem}
  \newtheorem{proposition}[proposition]{\notespropositionname}
  \aliascntresetthe{proposition}
  \newaliascnt{corollary}{theorem}
  \newtheorem{corollary}[corollary]{\notescorollaryname}
  \aliascntresetthe{corollary}

  \theoremstyle{definition}
  \newaliascnt{definition}{theorem}
  \newtheorem{definition}[definition]{\notesdefinitionname}
  \aliascntresetthe{definition}
  \newaliascnt{example}{theorem}
  \newtheorem{example}[example]{\notesexamplename}
  \aliascntresetthe{example}
  \newaliascnt{problem}{theorem}
  \newtheorem{problem}[problem]{\notesproblemname}
  \aliascntresetthe{problem}
  \newtheorem{exercise}{\notesexercisename}

  \theoremstyle{remark}
  \newaliascnt{remark}{theorem}
  \newtheorem{remark}[remark]{\notesremarkname}
  \aliascntresetthe{remark}

  \newenvironment{solution}
  {
  \begin{proof}[\notessolutionname]}
    {
  \end{proof}}

  \usepackage[nameinlink,noabbrev]{cleveref}

  \crefname{theorem}{theorem}{theorems}
  \crefname{lemma}{lemma}{lemmas}
  \crefname{proposition}{proposition}{propositions}
  \crefname{corollary}{corollary}{corollaries}
  \crefname{definition}{definition}{definitions}
  \crefname{example}{example}{examples}
  \crefname{problem}{problem}{problems}
  \crefname{exercise}{exercise}{exercises}
  \crefname{remark}{remark}{remarks}

  \ifx\notelanguage\notesfinnishlanguage
  \crefname{theorem}{lause}{lauseet}
  \crefname{lemma}{lemma}{lemmat}
  \crefname{proposition}{propositio}{propositiot}
  \crefname{corollary}{korollaari}{korollaarit}
  \crefname{definition}{määritelmä}{määritelmät}
  \crefname{example}{esimerkki}{esimerkit}
  \crefname{problem}{tehtävä}{tehtävät}
  \crefname{exercise}{tehtävä}{tehtävät}
  \crefname{remark}{huomautus}{huomautukset}
  \fi

  \newcommand{\NN}{\mathbb{N}}
  \newcommand{\NNone}{\mathbb{N}_1}
  \newcommand{\ZZ}{\mathbb{Z}}
  \newcommand{\QQ}{\mathbb{Q}}
  \newcommand{\RR}{\mathbb{R}}
  \newcommand{\CC}{\mathbb{C}}
  \newcommand{\PP}{\mathbb{P}}

  \DeclareMathOperator{\im}{im}
  \DeclareMathOperator{\rank}{rank}
  \DeclareMathOperator{\syt}{syt}
  \DeclareMathOperator{\pyj}{pyj}
  \DeclareMathOperator{\dom}{dom}
  \DeclarePairedDelimiter{\abs}{\lvert}{\rvert}
  \DeclarePairedDelimiter{\norm}{\lVert}{\rVert}
  \DeclarePairedDelimiterXPP{\maxnorm}[1]{}{\lVert}{\rVert}{_\infty}{#1}
  \newcommand{\set}[1]{\left\lbrace#1\right\rbrace}
  \newcommand{\maxset}[1]{\max\set{#1}}
  \newcommand{\minset}[1]{\min\set{#1}}
  \newcommand{\powerset}[1]{\mathcal{P}\!\left(#1\right)}


\title{Lineaarialgebra ja matriisilaskenta I\\Harjoitus 4}
\author{}
\date{}

\begin{document}

\exerciseheading{Lineaarialgebra ja matriisilaskenta I}{Harjoitus 4}

\begin{exercise}
  Olkoot
  \[
    v_1 =
    \begin{bmatrix}1 \\ 2 \\ 1
    \end{bmatrix},\quad
    v_2 =
    \begin{bmatrix}2 \\ 9 \\ 0
    \end{bmatrix},\quad
    v_3 =
    \begin{bmatrix}3 \\ 3 \\ 4
    \end{bmatrix},\quad
    u =
    \begin{bmatrix}5 \\ -1 \\ 9
    \end{bmatrix},\quad
    w =
    \begin{bmatrix}-1 \\ 3 \\ 2
    \end{bmatrix}
    \in\RR^{3\times 1}.
  \]
  \begin{enumerate}[label=(\alph*), leftmargin=*]
    \item Osoita, että jono \((v_1,v_2,v_3)\) on avaruuden
      \(\RR^{3\times 1}\) kanta.
    \item Määritä vektorin \(u\) koordinaattivektori kannassa
      \((v_1,v_2,v_3)\).
    \item Määritä vektori \(v\in\RR^{3\times 1}\), jonka
      koordinaattivektori kannassa \((v_1,v_2,v_3)\) on \(w\).
  \end{enumerate}
\end{exercise}

\begin{solution}
  Merkitään \(A=[v_1\ v_2\ v_3]\). Muokataan liitetty matriisi
  \([A\mid 0\ u]\) supistettuun porrasmuotoon:
  \begingroup
  \setlength{\arraycolsep}{4pt}
  \renewcommand{\arraystretch}{1.1}
  \[
    \begin{alignedat}{2}
      [A\mid 0\ u] &={}
      \left[
        \begin{array}{rrr|rr}
          1 & 2 & 3 & 0 & 5 \\
          2 & 9 & 3 & 0 & -1 \\
          1 & 0 & 4 & 0 & 9
        \end{array}
      \right]
      \quad
      \begin{array}{@{}l@{}}
        \smash{\rowop{R_1\leftrightarrow R_3}} \\
        {} \\
        {}
      \end{array}
      &\qquad\sim{}&
      \left[
        \begin{array}{rrr|rr}
          1 & 0 & 4 & 0 & 9 \\
          2 & 9 & 3 & 0 & -1 \\
          1 & 2 & 3 & 0 & 5
        \end{array}
      \right]
      \quad
      \begin{array}{@{}l@{}}
        {} \\
        \smash{\rowop{-2R_1}} \\
        \smash{\rowop{-R_1}}
      \end{array}
      \\[1.5ex]
      &\sim{}
      \left[
        \begin{array}{rrr|rr}
          1 & 0 & 4 & 0 & 9 \\
          0 & 9 & -5 & 0 & -19 \\
          0 & 2 & -1 & 0 & -4
        \end{array}
      \right]
      \quad
      \begin{array}{@{}l@{}}
        {} \\
        \smash{\rowop{\cdot 2}} \\
        {}
      \end{array}
      &\qquad\sim{}&
      \left[
        \begin{array}{rrr|rr}
          1 & 0 & 4 & 0 & 9 \\
          0 & 18 & -10 & 0 & -38 \\
          0 & 2 & -1 & 0 & -4
        \end{array}
      \right]
      \quad
      \begin{array}{@{}l@{}}
        {} \\
        \smash{\rowop{-9R_3}} \\
        {}
      \end{array}
      \\[1.5ex]
      &\sim{}
      \left[
        \begin{array}{rrr|rr}
          1 & 0 & 4 & 0 & 9 \\
          0 & 0 & -1 & 0 & -2 \\
          0 & 2 & -1 & 0 & -4
        \end{array}
      \right]
      \quad
      \begin{array}{@{}l@{}}
        \smash{\rowop{+4R_2}} \\
        {} \\
        \smash{\rowop{-R_2}}
      \end{array}
      &\qquad\sim{}&
      \left[
        \begin{array}{rrr|rr}
          1 & 0 & 0 & 0 & 1 \\
          0 & 0 & -1 & 0 & -2 \\
          0 & 2 & 0 & 0 & -2
        \end{array}
      \right]
      \quad
      \begin{array}{@{}l@{}}
        {} \\
        \smash{\rowop{\cdot(-1)}} \\
        \smash{\rowop{\cdot(1/2)}}
      \end{array}
      \\[1.5ex]
      &\sim{}
      \left[
        \begin{array}{rrr|rr}
          1 & 0 & 0 & 0 & 1 \\
          0 & 0 & 1 & 0 & 2 \\
          0 & 1 & 0 & 0 & -1
        \end{array}
      \right]
      \quad
      \begin{array}{@{}l@{}}
        {} \\
        \smash{\rowop{R_2\leftrightarrow R_3}} \\
        {}
      \end{array}
      &\qquad\sim{}&
      \left[
        \begin{array}{rrr|rr}
          1 & 0 & 0 & 0 & 1 \\
          0 & 1 & 0 & 0 & -1 \\
          0 & 0 & 1 & 0 & 2
        \end{array}
      \right].
    \end{alignedat}
  \]
  \endgroup
  \begin{enumerate}[label=(\alph*), leftmargin=*]
    \item Nollasaraketta vastaavasta yhtälöstä nähdään, että
      \(Ax=0\) toteutuu vain, kun \(x=0\). Lemman 3.6.3 nojalla
      jono \((v_1,v_2,v_3)\) on siis vapaa. Korollaarin 3.10.3
      perusteella se on avaruuden \(\RR^{3\times 1}\) kanta.

    \item Viimeisestä sarakkeesta saadaan yhtälön \(Ac=u\) ratkaisu
      \(c_1=1\), \(c_2=-1\) ja \(c_3=2\), eli \(u=v_1-v_2+2v_3\).
      Määritelmän 3.7.4 mukaan koordinaattivektori on siis
      \[
        c=
        \begin{bmatrix}1 \\ -1 \\ 2\end{bmatrix}.
      \]
      Kertoimet ovat yksikäsitteiset Lemman 3.7.3 perusteella.

    \item Koska vektorin \(v\) koordinaattivektori kannassa
      \((v_1,v_2,v_3)\) on \(w\), määritelmän 3.7.4 mukaan
      \[
        \begin{aligned}
          v &=-v_1+3v_2+2v_3=Aw \\
            &=\begin{bmatrix}
              1 & 2 & 3 \\
              2 & 9 & 3 \\
              1 & 0 & 4
            \end{bmatrix}
            \begin{bmatrix}-1 \\ 3 \\ 2\end{bmatrix}
            =\begin{bmatrix}11 \\ 31 \\ 7\end{bmatrix}.
        \end{aligned}
        \qedhere
      \]
  \end{enumerate}
\end{solution}

\newpage

\begin{exercise}
  Määritä matriisin
  \[
    A =
    \begin{bmatrix}
      1 & 4 & 5 & 2 \\
      2 & 1 & 3 & 0 \\
      -1 & 3 & 2 & 2
    \end{bmatrix}
  \]
  nolla-avaruuden kanta. Mikä on matriisin \(A\) aste?
\end{exercise}

\begin{exercise}
  Määritä matriisin
  \[
    A =
    \begin{bmatrix}
      1 & -3 & 4 & -2 & 5 & 4 \\
      2 & -6 & 9 & -1 & 8 & 2 \\
      2 & -6 & 9 & -1 & 9 & 7 \\
      -1 & 3 & -4 & 2 & -5 & -4
    \end{bmatrix}
  \]
  sarakeavaruuden kanta. Mikä on matriisin \(A\) aste?
\end{exercise}

\begin{exercise}
  Olkoot \(v=
    \begin{bmatrix}2 \\ k
  \end{bmatrix}\in\RR^{2\times 1}\)
  ja \(w=
    \begin{bmatrix}3 \\ 5
  \end{bmatrix}\in\RR^{2\times 1}\).
  Määritä \(k\in\RR\) siten, että
  \begin{enumerate}[label=(\alph*), leftmargin=*]
    \item \(v\) ja \(w\) ovat yhdensuuntaiset,
    \item \(v\) ja \(w\) ovat kohtisuorassa toisiaan vastaan,
    \item \(v\):n ja \(w\):n välinen kulma on \(\frac{\pi}{3}\).
  \end{enumerate}
  Piirrä kuva jokaisesta tapauksesta (a), (b) ja (c).
\end{exercise}

\begin{exercise}
  Olkoot
  \[
    v_1 =
    \begin{bmatrix}1 \\ 0 \\ 1 \\ 0
    \end{bmatrix},\quad
    v_2 =
    \begin{bmatrix}1 \\ 1 \\ 3 \\ 0
    \end{bmatrix},\quad
    v_3 =
    \begin{bmatrix}1 \\ 2 \\ 0 \\ 0
    \end{bmatrix}
    \in\RR^{4\times 1},
  \]
  ja olkoon \(V\) avaruuden \(\RR^{4\times 1}\) aliavaruus, jonka
  virittää jono \((v_1,v_2,v_3)\), siis
  \(V=\operatorname{Sp}(v_1,v_2,v_3)\).
  \begin{enumerate}[label=(\alph*), leftmargin=*]
    \item Osoita, että jono \((v_1,v_2,v_3)\) on vapaa, jolloin
      \((v_1,v_2,v_3)\) on \(V\):n kanta.
    \item Määritä \(V\):n ortonormaali kanta \((u_1,u_2,u_3)\)
      soveltamalla Gramin--Schmidtin menetelmää jonoon
      \((v_1,v_2,v_3)\).
  \end{enumerate}
\end{exercise}

\begin{exercise}
  Jatkoa edelliseen tehtävään. Täydennä ortonormaali jono
  \((u_1,u_2,u_3)\) avaruuden \(\RR^{4\times 1}\) ortonormaaliksi
  kannaksi.
\end{exercise}

\end{document}
